% Figure: geometry for the fatigue-crack-growth case study. A pressure-
% vessel wall strip carries a cyclic membrane (hoop) stress. A semi-
% elliptical surface flaw of depth a sits at the inner wall and grows
% through the thickness B under the cyclic stress. The leak-before-break
% condition is that the crack penetrates the wall (a = B, leak) before it
% reaches the critical depth at which fast fracture occurs (a = a_c, break).
\begin{figure}[htbp]
\centering
\begin{tikzpicture}[scale=1.0]
% wall strip (a slice of the vessel wall, thickness B horizontal)
\draw[engblack, very thick] (0,-1.6) rectangle (3.0,1.6);
% inner-wall face label (left edge)
\node[enggray, font=\scriptsize, anchor=east, rotate=90] at (-0.15,0) {inner wall};
\node[enggray, font=\scriptsize, anchor=west, rotate=90] at (3.2,0) {outer wall};
% wall thickness annotation
\draw[enggray, {Latex}-{Latex}] (0,-2.0) -- (3.0,-2.0);
\node[enggray, font=\scriptsize, anchor=north] at (1.5,-2.0) {wall thickness $B = 18\ \text{mm}$};
% the surface crack: a semi-ellipse on the inner wall growing rightward
\draw[engred, very thick, fill=engred!12]
  (0,0.55) .. controls (1.1,0.55) and (1.1,-0.55) .. (0,-0.55);
\node[engred, font=\scriptsize, anchor=west] at (1.15,0) {flaw, depth $a$};
% current crack depth annotation
\draw[engred, {Latex}-{Latex}] (0,0.9) -- (1.05,0.9);
\node[engred, font=\scriptsize, anchor=south] at (0.5,0.9) {$a_0 = 1.5\ \text{mm}$};
% cyclic hoop stress arrows (membrane tension across the strip, top and bottom)
\draw[engorange, very thick, -{Latex}] (0.6,1.65) -- (0.6,2.4);
\draw[engorange, very thick, -{Latex}] (1.8,1.65) -- (1.8,2.4);
\draw[engorange, very thick, -{Latex}] (2.6,1.65) -- (2.6,2.4);
\node[engorange, font=\small, anchor=south] at (1.6,2.4)
  {cyclic hoop stress $\Delta\sigma$};
\draw[engorange, very thick, {Latex}-] (0.6,-2.4) -- (0.6,-1.65);
\draw[engorange, very thick, {Latex}-] (1.8,-2.4) -- (1.8,-1.65);
\draw[engorange, very thick, {Latex}-] (2.6,-2.4) -- (2.6,-1.65);
% the two endpoints of the life: leak (a=B) and break (a=a_c)
\draw[englime!60!black, thick, dashed] (3.0,-1.6) -- (3.0,1.6);
\node[englime!60!black, font=\scriptsize, anchor=south, rotate=90]
  at (3.0,0.0) {};
\end{tikzpicture}
\caption[Surface flaw growing through a pressure-vessel wall]{Geometry
for the fatigue-crack-growth case study. A semi-elliptical surface flaw
of initial depth $a_0 = 1.5\ \text{mm}$ sits on the inner wall of a
pressure vessel and grows through the wall thickness $B = 18\ \text{mm}$
under the cyclic hoop stress range $\Delta\sigma$ delivered once per
pressurisation cycle. Two outcomes compete. If the crack reaches the
critical depth $a_c$ at which the stress-intensity factor equals the
fracture toughness, the wall fails by fast fracture (break). If instead
the crack penetrates the full wall while $a$ is still below $a_c$, the
vessel leaks at a detectable rate before it breaks. The design target is
the second outcome, leak-before-break.}
\label{fig:vol03ch03:edge-crack-geometry}
\end{figure}
